% BANKS-SOURCE: pdf=258 print=248 kind=appendix id=appendixC
\BanksAppendix{Diracology}

% BANKS-SOURCE: pdf=258 print=248 kind=prose
Here we collect a variety of identities for Dirac matrices and spinors. The
$4\times4$ matrices, $\gamma_5$, and Levi--Civita identities below are
four-dimensional ($D=4$). We use
$\eta_{\mu\nu}=\operatorname{diag}(-,+,+,+)$ and
$\epsilon^{0123}=+1$, so $\epsilon_{0123}=-1$. We write
$\slashed p=\gamma^\mu p_\mu$. The basic commutation relations are

% BANKS-SOURCE: pdf=258 print=248 kind=equation id=C.1
\begin{equation}
[\gamma^\mu,\gamma^\nu]_+=-2\eta^{\mu\nu}.
\label{eq:C.1}
\end{equation}

The Weyl representation of these relations is

% BANKS-SOURCE: pdf=258 print=248 kind=equation id=C.2
\begin{equation}
\gamma^\mu=\begin{pmatrix}0&\sigma^\mu\\[2pt]\bar\sigma^\mu&0\end{pmatrix}.
\label{eq:C.2}
\end{equation}

where $\sigma^\mu\equiv(1,\boldsymbol{\sigma})$,
$\bar\sigma^\mu\equiv(1,-\boldsymbol{\sigma})$, and $\boldsymbol{\sigma}$ is
the usual 3-vector of Pauli matrices. In this representation

% BANKS-SOURCE: pdf=258 print=248 kind=display
\[
\gamma_5\equiv\ii\gamma^0\gamma^1\gamma^2\gamma^3
=\begin{pmatrix}-1&0\\0&1\end{pmatrix}.
\]
In $D=4$, $\gamma_5^2=1$ and
$\{\gamma_5,\gamma^\mu\}=0$, so the chirality projectors are
$P_L=(1-\gamma_5)/2$ and $P_R=(1+\gamma_5)/2$.

These matrices obviously satisfy

% BANKS-SOURCE: pdf=258 print=248 kind=equation id=C.3
\begin{equation}
(\gamma^\mu)^\dagger=\gamma^0\gamma^\mu\gamma^0.
\label{eq:C.3}
\end{equation}
\BanksImplicitHook{I-APPC-001}

We will work only with representations related to the Weyl representation by unitary
(rather than special linear) transformations, so this relation will always be true.

A convenient representation is the Dirac representation, in which

% BANKS-SOURCE: pdf=258 print=248 kind=display
\[
\gamma_D^0=\begin{pmatrix}\mathbf 1&0\\0&-\mathbf 1\end{pmatrix}.
\]

It is related to the Weyl representation by $\gamma_D^\mu=S_D^\dagger\gamma_W^\mu S_D$ with

% BANKS-SOURCE: pdf=258 print=248 kind=display
\[
S_D=\frac{1}{\sqrt2}(1-\ii\sigma_2)\otimes1.
\]

The Majorana representation is related to the Dirac representation by $\gamma_M^\mu=S_M^\dagger\gamma_D^\mu S_M$
with

% BANKS-SOURCE: pdf=258 print=248 kind=display
\[
S_M=\frac{1}{\sqrt2}
\begin{pmatrix}
\mathbf 1&\sigma_2\\
\sigma_2&-\mathbf 1
\end{pmatrix}.
\]

In the Majorana representation the Dirac equation is satisfied separately by the real
and imaginary parts of the field.

% BANKS-SOURCE: pdf=259 print=249 kind=prose
The space of all $4\times4$ matrices is spanned by the anti-symmetrized products of Dirac
matrices, $\gamma^{\mu_1\ldots\mu_k}\equiv
\gamma^{[\mu_1}\cdots\gamma^{\mu_k]}$, with $0\leq k\leq4$. In $D=4$, the
Hodge-duality relations are

% BANKS-SOURCE: pdf=259 print=249 kind=display
\[
\gamma^{\mu\nu\lambda\kappa}=-\ii\epsilon^{\mu\nu\lambda\kappa}\gamma_5,
\qquad
\gamma^{\mu\nu\lambda}=-\ii\epsilon^{\mu\nu\lambda\kappa}\gamma_\kappa\gamma_5,
\]

and

% BANKS-SOURCE: pdf=259 print=249 kind=display
\[
\gamma^{\mu\nu}=\frac{\ii}{2}\epsilon^{\mu\nu\lambda\kappa}
\gamma_{\lambda\kappa}\gamma_5.
\]

Calculations of spin-averaged cross sections and closed fermion loops lead to traces
of products of Dirac matrices:

% BANKS-SOURCE: pdf=259 print=249 kind=display
\[
\operatorname{tr}(\gamma^{\mu_1}\ldots\gamma^{\mu_k}).
\]

These can be calculated using anti-commutation relations and cyclicity of the trace, or
by the tensor method. The latter comes from the observation that the result must be a
numerical Lorentz tensor (and thus built from $\eta_{\mu\nu}$ and $\epsilon_{\mu\nu\lambda\kappa}$), which is invariant under
cyclic permutation of the indices. Note also that the result is zero for $k$ odd, because
$\gamma_5^2=1$ and $\gamma_5$ anti-commutes with all the $\gamma^\mu$. For example

% BANKS-SOURCE: pdf=259 print=249 kind=display
\[
\operatorname{tr}(\gamma^\mu\gamma^\nu)=A\eta^{\mu\nu},
\]

and we compute $A=-4$ by hitting both sides with $\eta_{\mu\nu}$. A similar
manipulation shows that

% BANKS-SOURCE: pdf=259 print=249 kind=display
\[
\operatorname{tr}(\gamma^\mu\gamma^\nu\gamma_5)=0,
\]

% BANKS-SOURCE: pdf=259 print=249 kind=display
\[
\operatorname{tr}(\gamma^\mu\gamma^\nu\gamma^\lambda\gamma^\kappa)
=C_1(\eta^{\mu\nu}\eta^{\lambda\kappa}+\eta^{\mu\kappa}\eta^{\lambda\nu})
+C_2\eta^{\mu\lambda}\eta^{\nu\kappa}+B\epsilon^{\mu\nu\lambda\kappa}.
\]

Taking all the indices different, the RHS is just $B$ and the LHS is $\propto\operatorname{tr}\gamma_5=0$, so $B=0$.
On hitting both sides with $\eta_{\mu\nu}$, we get

% BANKS-SOURCE: pdf=259 print=249 kind=display
\[
5C_1+C_2=16\qquad(D=4).
\]

Doing the same with $\eta_{\mu\lambda}$ gives

% BANKS-SOURCE: pdf=259 print=249 kind=display
\[
\operatorname{tr}(\gamma^\mu\gamma^\nu\gamma_\mu\gamma^\kappa)
=-8\eta^{\nu\kappa}=(2C_1+4C_2)\eta^{\nu\kappa}
\qquad(D=4).
\]

Here we have used an identity

% BANKS-SOURCE: pdf=259 print=249 kind=display
\[
\gamma^\mu\gamma^\nu\gamma_\mu=2\gamma^\nu\qquad(D=4),
\]

which we will derive in a moment. We conclude that $C_1=-C_2=4$.
Similar manipulations allow us to calculate any trace with relative ease.

% BANKS-SOURCE: pdf=260 print=250 kind=prose
To evaluate $\gamma^\lambda\gamma_{\mu_1}\ldots\gamma_{\mu_k}\gamma_\lambda$, we
proceed by induction, using the anti-commutation relations. The first
contraction is valid in $D$ spacetime dimensions:

% BANKS-SOURCE: pdf=260 print=250 kind=display
\[
\gamma^\lambda\gamma_\mu\gamma_\lambda
=D\gamma_\mu-2\gamma_\mu=(D-2)\gamma_\mu.
\]

% BANKS-SOURCE: pdf=260 print=250 kind=display
\[
\gamma^\lambda\gamma_\mu\gamma_\nu\gamma_\lambda
=-\gamma_\mu(2\gamma_\nu)-2\gamma_\nu\gamma_\mu
=4\eta_{\mu\nu}\qquad(D=4),
\]

% BANKS-SOURCE: pdf=260 print=250 kind=display
\[
\gamma^\lambda\gamma_\mu\gamma_\nu\gamma_\kappa\gamma_\lambda
=-4\gamma_\mu\eta_{\nu\kappa}-2\gamma_\nu\gamma_\kappa\gamma_\mu
\qquad(D=4),
\]

% BANKS-SOURCE: pdf=260 print=250 kind=display
\[
\gamma^\lambda\gamma_\mu\gamma_\nu\gamma_\kappa\gamma_\alpha\gamma_\lambda
=4\gamma_\mu\gamma_\nu\eta_{\kappa\alpha}
+2\gamma_\mu\gamma_\kappa\gamma_\alpha\gamma_\nu
-2\gamma_\nu\gamma_\kappa\gamma_\alpha\gamma_\mu
\qquad(D=4).
\]

The reader is encouraged to continue computing these trace and contraction identities
to higher orders. The $D$ in the first line is the physical spacetime dimension and is
independent of the dimensional-regularization parameter $d_{\mathrm{reg}}$.
\BanksImplicitHook{I-APPC-002}

The wave functions which convert fermion creation operators into local fields are
solutions of the momentum-space Dirac equations

With positive-frequency modes $\exp(+\ii p\mathbin{\cdot}x)$, where
$p\mathbin{\cdot}x=-Et+\mathbf p\mathbin{\cdot}\mathbf x$, the kinetic operator
$-\ii\slashed\partial-m$ gives

% BANKS-SOURCE: pdf=260 print=250 kind=display
\[
(\slashed p-m)u(p,s)=(\slashed p+m)v(p,s)=0.
\]

They satisfy the completeness relations

% BANKS-SOURCE: pdf=260 print=250 kind=display
\[
\sum_su(p,s)\bar u(p,s)=-(\slashed p+m),
\]

% BANKS-SOURCE: pdf=260 print=250 kind=display
\[
\sum_sv(p,s)\bar v(p,s)=-(\slashed p-m),
\]

where the right-hand sides are proportional to the projection operators on solutions
of the respective equations. The minus signs follow from the mostly-plus
Clifford convention and the Dirac adjoint $\bar u=u^\dagger\gamma^0$. Note that
many books use a different normalization. We have followed the convention that,
for fields of any spin in $D$ spacetime dimensions, the Fourier transform is
$(2E_{\mathbf p}(2\pi)^{D-1})^{-1/2}$ times a creation or annihilation operator,
times a covariant wave function.

In the Weyl basis, the solutions have the form

% BANKS-SOURCE: pdf=260 print=250 kind=display
\[
u(p,s)=\begin{pmatrix}\sqrt{-p\cdot\sigma}\,\chi(s)\\[2pt]-\sqrt{-p\cdot\bar\sigma}\,\chi(s)\end{pmatrix},
\qquad
v(p,s)=\begin{pmatrix}\sqrt{-p\cdot\sigma}\,\eta(s)\\[2pt]\sqrt{-p\cdot\bar\sigma}\,\eta(s)\end{pmatrix}.
\]

Here $p^\mu=(E_{\mathbf p},\mathbf p)$, $p_\mu=(-E_{\mathbf p},\mathbf p)$,
and $p^2=-m^2$, so $p\cdot\sigma=p_\mu\sigma^\mu$ and
$p\cdot\bar\sigma=p_\mu\bar\sigma^\mu$; the square roots above are
positive for $E_{\mathbf p}>0$. The formulas are the $D=4$ Weyl-basis
specialization. The spinors $\chi$ and $\eta$ are normalized two-component
spinors corresponding to the choice of $s$.
In a convenient convention

% BANKS-SOURCE: pdf=260 print=250 kind=display
\[
\eta(s)=2s\ii\sigma_2\chi^*(s),
\]

where $s=\pm1/2$ is either the component of spin in some fixed direction, $\widehat{\mathbf n}$, or the
helicity.
